% =============================================================================
% Assignments/CS301/03-arrays-stacks-coding-assignment.tex
% Coding Assignment 3: Stack over Array + Bracket Check + Postfix Evaluation
%
% Student-facing problem statement. Instructor-only files live in
% Assignments/CS301/03-arrays-stacks-coding/:
%   reference_solution.c  (NEVER published)
%   test_hidden.c         (NEVER published)
%   autograder.sh         (NEVER published)
%   stack.h, starter_stack.c, test_visible.c
%
% Compile with: cd Assignments/CS301 &&
%   TEXINPUTS=../../Preambles;$TEXINPUTS xelatex -interaction=nonstopmode 03-arrays-stacks-coding-assignment.tex
%   (Windows/MiKTeX uses ; not : as the path separator)
% =============================================================================

\documentclass[11pt]{article}
\usepackage[margin=1in]{geometry}
% =============================================================================
% Preambles/header.tex — shared LaTeX/Beamer preamble for this project
%
% Use from a Beamer lecture by prepending:
%     % =============================================================================
% Preambles/header.tex — shared LaTeX/Beamer preamble for this project
%
% Use from a Beamer lecture by prepending:
%     % =============================================================================
% Preambles/header.tex — shared LaTeX/Beamer preamble for this project
%
% Use from a Beamer lecture by prepending:
%     \input{header}
% and compiling with TEXINPUTS=../Preambles:$TEXINPUTS (the /compile-latex
% skill does this automatically).
%
% PALETTE CONTRACT. Color names defined here MUST match the names in
% Quarto/theme-template.scss so Beamer and Quarto renderings use the same
% palette. The scripts/check-palette-sync.sh script enforces this.
%
% To customize: edit the HEX values below AND the matching $variable in
% Quarto/theme-template.scss, then run scripts/check-palette-sync.sh.
%
% TITLE PAGE. A formal, serif (Libertine) title-page template is defined in
% the Beamer-only block below (\setbeamertemplate{title page}) -- title,
% subtitle, a thin gold rule, then \author/\institute/\date. Frametitle and
% body text remain on Beamer's sans-serif default. No SCSS counterpart is
% needed for this (font/layout only, no new colors).
% =============================================================================

% -----------------------------------------------------------------------------
% Engine detection + encoding
%
% Our /compile-latex skill uses XeLaTeX, where inputenc is unnecessary and
% can produce encoding warnings. Only load inputenc under pdfTeX.
% -----------------------------------------------------------------------------
\usepackage{iftex}
\ifPDFTeX
  \usepackage[utf8]{inputenc}
\fi

\usepackage{amsmath, amssymb, amsthm}
\usepackage{booktabs}
\usepackage{graphicx}

% -----------------------------------------------------------------------------
% Serif face for the title page (formal/academic look). Linux Libertine is
% XeLaTeX- and pdfTeX-safe and redefines \rmdefault, so \rmfamily picks it up
% everywhere \rmfamily is used explicitly. Body text and frametitles stay on
% Beamer's own sans-serif default (unaffected) -- only the title-page
% \setbeamerfont{...} overrides below opt in to \rmfamily.
% -----------------------------------------------------------------------------
\usepackage{libertine}

% -----------------------------------------------------------------------------
% PALETTE — mirrors Quarto/theme-template.scss variable names.
% Load xcolor and define colors BEFORE anything that references them
% (notably hyperref, hypersetup, and the Beamer color assignments).
% -----------------------------------------------------------------------------
\usepackage{xcolor}

\definecolor{primary-blue}{HTML}{012169}      % headings, accents
\definecolor{primary-gold}{HTML}{B9975B}      % emphasis, borders
\definecolor{highlight-yellow}{HTML}{F2A900}  % markers, alerts
\definecolor{light-bg}{HTML}{E8EDF5}          % title / section backgrounds
\definecolor{jet}{HTML}{1A1A1A}               % body text

% Semantic colors (used in diagrams and annotations)
\definecolor{positive}{HTML}{15803D}          % good / observed / identified
\definecolor{negative}{HTML}{B91C1C}          % bad / problematic / bias
\definecolor{neutral}{HTML}{525252}           % reference / context

% Highlight accents (match .hi-* classes in the SCSS)
\definecolor{hi-slate}{HTML}{314F4F}
\definecolor{hi-green}{HTML}{2E8B57}
\definecolor{hi-red}{HTML}{C41E3A}

% Semantic aliases so skill templates and snippets can write
% \textcolor{accent}{...} without hard-coding "primary-blue".
\colorlet{accent}{primary-blue}
\colorlet{accent2}{primary-gold}

% -----------------------------------------------------------------------------
% Hyperref — loaded AFTER xcolor + palette so link colors resolve.
% -----------------------------------------------------------------------------
\usepackage{hyperref}
\hypersetup{colorlinks=true, linkcolor=primary-blue, urlcolor=primary-blue,
            citecolor=primary-blue, pdfborder={0 0 0}}

% -----------------------------------------------------------------------------
% Course code — set once per deck (after \input{header}, before \begin{document})
% via \coursecode{CS401}. Shown in the footer on every slide so a deck's course
% is identifiable without opening the title page. Empty by default.
% -----------------------------------------------------------------------------
\makeatletter
\newcommand{\coursecode}[1]{\gdef\@coursecodetext{#1}}
\gdef\@coursecodetext{}
\makeatother

% -----------------------------------------------------------------------------
% Beamer theme assignments (only apply under Beamer).
%
% We detect Beamer via \@ifclassloaded{beamer}, which is the documented,
% non-internal way. This must be wrapped in \makeatletter/\makeatother
% because \@ifclassloaded contains an @-catcode character.
% -----------------------------------------------------------------------------
\makeatletter
\@ifclassloaded{beamer}{%
  \usefonttheme{professionalfonts}
  \setbeamercolor{structure}{fg=primary-blue}
  \setbeamercolor{frametitle}{fg=primary-blue}
  \setbeamercolor{title}{fg=primary-blue}
  \setbeamercolor{subtitle}{fg=primary-gold}
  \setbeamercolor{author}{fg=primary-blue}
  \setbeamercolor{institute}{fg=neutral}
  \setbeamercolor{date}{fg=primary-gold}
  \setbeamercolor{itemize item}{fg=primary-blue}
  \setbeamercolor{itemize subitem}{fg=primary-gold}
  \setbeamercolor{alerted text}{fg=primary-gold}
  \setbeamercolor{block title}{fg=white, bg=primary-blue}
  \setbeamercolor{block body}{bg=light-bg}
  % Semantic block variants: block=definition/concept (blue), exampleblock=
  % worked example (green/positive), alertblock=key takeaway or caution (gold).
  % Keep to <=2 colored boxes per slide across ALL three variants (INV-7).
  \setbeamercolor{block title example}{fg=white, bg=positive}
  \setbeamercolor{block body example}{bg=light-bg}
  \setbeamercolor{block title alerted}{fg=white, bg=primary-gold}
  \setbeamercolor{block body alerted}{bg=light-bg}

  % ---------------------------------------------------------------------------
  % Formal title page: serif (Libertine, via \rmfamily) for title-page
  % elements only. Body text, itemize, and frametitle are intentionally left
  % on Beamer's sans-serif default -- \usebeamerfont{title} is also what
  % \transitionslide (below) uses for its big centered text, so section
  % transitions pick up the same serif treatment for free.
  % ---------------------------------------------------------------------------
  \setbeamerfont{title}{family=\rmfamily, series=\bfseries, size=\Large}
  \setbeamerfont{subtitle}{family=\rmfamily, shape=\itshape, size=\normalsize}
  \setbeamerfont{author}{family=\rmfamily, size=\normalsize}
  \setbeamerfont{institute}{family=\rmfamily, size=\scriptsize}
  \setbeamerfont{date}{family=\rmfamily, size=\scriptsize}

  \setbeamertemplate{title page}{
    \vfill
    \begin{center}
      {\usebeamerfont{title}\usebeamercolor[fg]{title}\inserttitle\par}
      \vspace{0.15cm}
      {\usebeamerfont{subtitle}\usebeamercolor[fg]{subtitle}\insertsubtitle\par}
      \vspace{0.45cm}
      {\color{primary-gold}\rule{0.42\textwidth}{0.6pt}\par}
      \vspace{0.45cm}
      {\usebeamerfont{author}\usebeamercolor[fg]{author}\insertauthor\par}
      \vspace{0.35cm}
      {\usebeamerfont{institute}\usebeamercolor[fg]{institute}\insertinstitute\par}
      \vspace{0.35cm}
      {\usebeamerfont{date}\usebeamercolor[fg]{date}\insertdate\par}
    \end{center}
    \vfill
  }

  % Minimal footer: course code (left) + page X / N (right), no navigation clutter
  \setbeamertemplate{navigation symbols}{}
  \setbeamertemplate{footline}{%
    \color{neutral}\footnotesize\hspace{0.7em}\@coursecodetext\hfill
    \insertframenumber/\inserttotalframenumber\hspace{0.7em}\vspace{0.3em}}
}{}
\makeatother

% -----------------------------------------------------------------------------
% TikZ setup — libraries the snippet gallery uses
% -----------------------------------------------------------------------------
\usepackage{tikz}
\usetikzlibrary{arrows.meta, positioning, calc, decorations.pathreplacing,
                fit, shapes.geometric, backgrounds}

% Shared TikZ styles — snippets and hand-written diagrams can pick these up
% via `\begin{tikzpicture}[every node/.style={...}]` or by naming specific
% styles (e.g., `\node[dag-node] (X) at (0,0) {X};`).
\tikzset{
  % Node styles for causal DAGs and similar
  dag-node/.style = {
    draw=primary-blue, fill=light-bg, rounded corners=2pt,
    minimum width=1.2cm, minimum height=0.9cm, align=center, font=\small
  },
  decision-node/.style = {
    draw=primary-gold, fill=white, diamond, aspect=1.8,
    minimum width=1.8cm, minimum height=1.2cm, align=center, font=\small,
    inner sep=1pt
  },
  % Edge styles — observed vs counterfactual
  observed-edge/.style = {-{Latex[length=2.5mm]}, thick, draw=jet},
  counterfactual-edge/.style = {-{Latex[length=2.5mm]}, thick, draw=neutral, dashed},
  confound-edge/.style = {-{Latex[length=2.5mm]}, thick, draw=negative, dashed},
  % Data-point styles for plots
  observed-dot/.style = {fill=primary-blue, draw=primary-blue, circle, inner sep=1.2pt},
  counterfactual-dot/.style = {fill=white, draw=neutral, circle, inner sep=1.2pt}
}

% -----------------------------------------------------------------------------
% Convenience macros
% -----------------------------------------------------------------------------
\newcommand{\muted}[1]{\textcolor{neutral}{#1}}
\newcommand{\key}[1]{\textcolor{primary-gold}{\textbf{#1}}}
\newcommand{\good}[1]{\textcolor{positive}{#1}}
\newcommand{\bad}[1]{\textcolor{negative}{#1}}

% Standout frame macro: full-bleed dark background for section transitions.
% Beamer-only (uses \begin{frame}/\setbeamercolor/\usebeamerfont) -- do not
% call from an article-class file (e.g. Notes/<CODE>/*-notes.tex). \muted,
% \key, \good, \bad, and \coursecode above are plain xcolor/gdef and are
% safe to use from either class; verified when Notes/article-class support
% was added.
% Expand: \transitionslide{Your transition title}
\newcommand{\transitionslide}[1]{%
  \begingroup
  \setbeamercolor{background canvas}{bg=primary-blue}
  \begin{frame}[plain]
    \centering
    \vfill
    {\usebeamerfont{title}\color{white}\Large #1\par}
    \vfill
  \end{frame}
  \endgroup
}

% and compiling with TEXINPUTS=../Preambles:$TEXINPUTS (the /compile-latex
% skill does this automatically).
%
% PALETTE CONTRACT. Color names defined here MUST match the names in
% Quarto/theme-template.scss so Beamer and Quarto renderings use the same
% palette. The scripts/check-palette-sync.sh script enforces this.
%
% To customize: edit the HEX values below AND the matching $variable in
% Quarto/theme-template.scss, then run scripts/check-palette-sync.sh.
%
% TITLE PAGE. A formal, serif (Libertine) title-page template is defined in
% the Beamer-only block below (\setbeamertemplate{title page}) -- title,
% subtitle, a thin gold rule, then \author/\institute/\date. Frametitle and
% body text remain on Beamer's sans-serif default. No SCSS counterpart is
% needed for this (font/layout only, no new colors).
% =============================================================================

% -----------------------------------------------------------------------------
% Engine detection + encoding
%
% Our /compile-latex skill uses XeLaTeX, where inputenc is unnecessary and
% can produce encoding warnings. Only load inputenc under pdfTeX.
% -----------------------------------------------------------------------------
\usepackage{iftex}
\ifPDFTeX
  \usepackage[utf8]{inputenc}
\fi

\usepackage{amsmath, amssymb, amsthm}
\usepackage{booktabs}
\usepackage{graphicx}

% -----------------------------------------------------------------------------
% Serif face for the title page (formal/academic look). Linux Libertine is
% XeLaTeX- and pdfTeX-safe and redefines \rmdefault, so \rmfamily picks it up
% everywhere \rmfamily is used explicitly. Body text and frametitles stay on
% Beamer's own sans-serif default (unaffected) -- only the title-page
% \setbeamerfont{...} overrides below opt in to \rmfamily.
% -----------------------------------------------------------------------------
\usepackage{libertine}

% -----------------------------------------------------------------------------
% PALETTE — mirrors Quarto/theme-template.scss variable names.
% Load xcolor and define colors BEFORE anything that references them
% (notably hyperref, hypersetup, and the Beamer color assignments).
% -----------------------------------------------------------------------------
\usepackage{xcolor}

\definecolor{primary-blue}{HTML}{012169}      % headings, accents
\definecolor{primary-gold}{HTML}{B9975B}      % emphasis, borders
\definecolor{highlight-yellow}{HTML}{F2A900}  % markers, alerts
\definecolor{light-bg}{HTML}{E8EDF5}          % title / section backgrounds
\definecolor{jet}{HTML}{1A1A1A}               % body text

% Semantic colors (used in diagrams and annotations)
\definecolor{positive}{HTML}{15803D}          % good / observed / identified
\definecolor{negative}{HTML}{B91C1C}          % bad / problematic / bias
\definecolor{neutral}{HTML}{525252}           % reference / context

% Highlight accents (match .hi-* classes in the SCSS)
\definecolor{hi-slate}{HTML}{314F4F}
\definecolor{hi-green}{HTML}{2E8B57}
\definecolor{hi-red}{HTML}{C41E3A}

% Semantic aliases so skill templates and snippets can write
% \textcolor{accent}{...} without hard-coding "primary-blue".
\colorlet{accent}{primary-blue}
\colorlet{accent2}{primary-gold}

% -----------------------------------------------------------------------------
% Hyperref — loaded AFTER xcolor + palette so link colors resolve.
% -----------------------------------------------------------------------------
\usepackage{hyperref}
\hypersetup{colorlinks=true, linkcolor=primary-blue, urlcolor=primary-blue,
            citecolor=primary-blue, pdfborder={0 0 0}}

% -----------------------------------------------------------------------------
% Course code — set once per deck (after % =============================================================================
% Preambles/header.tex — shared LaTeX/Beamer preamble for this project
%
% Use from a Beamer lecture by prepending:
%     \input{header}
% and compiling with TEXINPUTS=../Preambles:$TEXINPUTS (the /compile-latex
% skill does this automatically).
%
% PALETTE CONTRACT. Color names defined here MUST match the names in
% Quarto/theme-template.scss so Beamer and Quarto renderings use the same
% palette. The scripts/check-palette-sync.sh script enforces this.
%
% To customize: edit the HEX values below AND the matching $variable in
% Quarto/theme-template.scss, then run scripts/check-palette-sync.sh.
%
% TITLE PAGE. A formal, serif (Libertine) title-page template is defined in
% the Beamer-only block below (\setbeamertemplate{title page}) -- title,
% subtitle, a thin gold rule, then \author/\institute/\date. Frametitle and
% body text remain on Beamer's sans-serif default. No SCSS counterpart is
% needed for this (font/layout only, no new colors).
% =============================================================================

% -----------------------------------------------------------------------------
% Engine detection + encoding
%
% Our /compile-latex skill uses XeLaTeX, where inputenc is unnecessary and
% can produce encoding warnings. Only load inputenc under pdfTeX.
% -----------------------------------------------------------------------------
\usepackage{iftex}
\ifPDFTeX
  \usepackage[utf8]{inputenc}
\fi

\usepackage{amsmath, amssymb, amsthm}
\usepackage{booktabs}
\usepackage{graphicx}

% -----------------------------------------------------------------------------
% Serif face for the title page (formal/academic look). Linux Libertine is
% XeLaTeX- and pdfTeX-safe and redefines \rmdefault, so \rmfamily picks it up
% everywhere \rmfamily is used explicitly. Body text and frametitles stay on
% Beamer's own sans-serif default (unaffected) -- only the title-page
% \setbeamerfont{...} overrides below opt in to \rmfamily.
% -----------------------------------------------------------------------------
\usepackage{libertine}

% -----------------------------------------------------------------------------
% PALETTE — mirrors Quarto/theme-template.scss variable names.
% Load xcolor and define colors BEFORE anything that references them
% (notably hyperref, hypersetup, and the Beamer color assignments).
% -----------------------------------------------------------------------------
\usepackage{xcolor}

\definecolor{primary-blue}{HTML}{012169}      % headings, accents
\definecolor{primary-gold}{HTML}{B9975B}      % emphasis, borders
\definecolor{highlight-yellow}{HTML}{F2A900}  % markers, alerts
\definecolor{light-bg}{HTML}{E8EDF5}          % title / section backgrounds
\definecolor{jet}{HTML}{1A1A1A}               % body text

% Semantic colors (used in diagrams and annotations)
\definecolor{positive}{HTML}{15803D}          % good / observed / identified
\definecolor{negative}{HTML}{B91C1C}          % bad / problematic / bias
\definecolor{neutral}{HTML}{525252}           % reference / context

% Highlight accents (match .hi-* classes in the SCSS)
\definecolor{hi-slate}{HTML}{314F4F}
\definecolor{hi-green}{HTML}{2E8B57}
\definecolor{hi-red}{HTML}{C41E3A}

% Semantic aliases so skill templates and snippets can write
% \textcolor{accent}{...} without hard-coding "primary-blue".
\colorlet{accent}{primary-blue}
\colorlet{accent2}{primary-gold}

% -----------------------------------------------------------------------------
% Hyperref — loaded AFTER xcolor + palette so link colors resolve.
% -----------------------------------------------------------------------------
\usepackage{hyperref}
\hypersetup{colorlinks=true, linkcolor=primary-blue, urlcolor=primary-blue,
            citecolor=primary-blue, pdfborder={0 0 0}}

% -----------------------------------------------------------------------------
% Course code — set once per deck (after \input{header}, before \begin{document})
% via \coursecode{CS401}. Shown in the footer on every slide so a deck's course
% is identifiable without opening the title page. Empty by default.
% -----------------------------------------------------------------------------
\makeatletter
\newcommand{\coursecode}[1]{\gdef\@coursecodetext{#1}}
\gdef\@coursecodetext{}
\makeatother

% -----------------------------------------------------------------------------
% Beamer theme assignments (only apply under Beamer).
%
% We detect Beamer via \@ifclassloaded{beamer}, which is the documented,
% non-internal way. This must be wrapped in \makeatletter/\makeatother
% because \@ifclassloaded contains an @-catcode character.
% -----------------------------------------------------------------------------
\makeatletter
\@ifclassloaded{beamer}{%
  \usefonttheme{professionalfonts}
  \setbeamercolor{structure}{fg=primary-blue}
  \setbeamercolor{frametitle}{fg=primary-blue}
  \setbeamercolor{title}{fg=primary-blue}
  \setbeamercolor{subtitle}{fg=primary-gold}
  \setbeamercolor{author}{fg=primary-blue}
  \setbeamercolor{institute}{fg=neutral}
  \setbeamercolor{date}{fg=primary-gold}
  \setbeamercolor{itemize item}{fg=primary-blue}
  \setbeamercolor{itemize subitem}{fg=primary-gold}
  \setbeamercolor{alerted text}{fg=primary-gold}
  \setbeamercolor{block title}{fg=white, bg=primary-blue}
  \setbeamercolor{block body}{bg=light-bg}
  % Semantic block variants: block=definition/concept (blue), exampleblock=
  % worked example (green/positive), alertblock=key takeaway or caution (gold).
  % Keep to <=2 colored boxes per slide across ALL three variants (INV-7).
  \setbeamercolor{block title example}{fg=white, bg=positive}
  \setbeamercolor{block body example}{bg=light-bg}
  \setbeamercolor{block title alerted}{fg=white, bg=primary-gold}
  \setbeamercolor{block body alerted}{bg=light-bg}

  % ---------------------------------------------------------------------------
  % Formal title page: serif (Libertine, via \rmfamily) for title-page
  % elements only. Body text, itemize, and frametitle are intentionally left
  % on Beamer's sans-serif default -- \usebeamerfont{title} is also what
  % \transitionslide (below) uses for its big centered text, so section
  % transitions pick up the same serif treatment for free.
  % ---------------------------------------------------------------------------
  \setbeamerfont{title}{family=\rmfamily, series=\bfseries, size=\Large}
  \setbeamerfont{subtitle}{family=\rmfamily, shape=\itshape, size=\normalsize}
  \setbeamerfont{author}{family=\rmfamily, size=\normalsize}
  \setbeamerfont{institute}{family=\rmfamily, size=\scriptsize}
  \setbeamerfont{date}{family=\rmfamily, size=\scriptsize}

  \setbeamertemplate{title page}{
    \vfill
    \begin{center}
      {\usebeamerfont{title}\usebeamercolor[fg]{title}\inserttitle\par}
      \vspace{0.15cm}
      {\usebeamerfont{subtitle}\usebeamercolor[fg]{subtitle}\insertsubtitle\par}
      \vspace{0.45cm}
      {\color{primary-gold}\rule{0.42\textwidth}{0.6pt}\par}
      \vspace{0.45cm}
      {\usebeamerfont{author}\usebeamercolor[fg]{author}\insertauthor\par}
      \vspace{0.35cm}
      {\usebeamerfont{institute}\usebeamercolor[fg]{institute}\insertinstitute\par}
      \vspace{0.35cm}
      {\usebeamerfont{date}\usebeamercolor[fg]{date}\insertdate\par}
    \end{center}
    \vfill
  }

  % Minimal footer: course code (left) + page X / N (right), no navigation clutter
  \setbeamertemplate{navigation symbols}{}
  \setbeamertemplate{footline}{%
    \color{neutral}\footnotesize\hspace{0.7em}\@coursecodetext\hfill
    \insertframenumber/\inserttotalframenumber\hspace{0.7em}\vspace{0.3em}}
}{}
\makeatother

% -----------------------------------------------------------------------------
% TikZ setup — libraries the snippet gallery uses
% -----------------------------------------------------------------------------
\usepackage{tikz}
\usetikzlibrary{arrows.meta, positioning, calc, decorations.pathreplacing,
                fit, shapes.geometric, backgrounds}

% Shared TikZ styles — snippets and hand-written diagrams can pick these up
% via `\begin{tikzpicture}[every node/.style={...}]` or by naming specific
% styles (e.g., `\node[dag-node] (X) at (0,0) {X};`).
\tikzset{
  % Node styles for causal DAGs and similar
  dag-node/.style = {
    draw=primary-blue, fill=light-bg, rounded corners=2pt,
    minimum width=1.2cm, minimum height=0.9cm, align=center, font=\small
  },
  decision-node/.style = {
    draw=primary-gold, fill=white, diamond, aspect=1.8,
    minimum width=1.8cm, minimum height=1.2cm, align=center, font=\small,
    inner sep=1pt
  },
  % Edge styles — observed vs counterfactual
  observed-edge/.style = {-{Latex[length=2.5mm]}, thick, draw=jet},
  counterfactual-edge/.style = {-{Latex[length=2.5mm]}, thick, draw=neutral, dashed},
  confound-edge/.style = {-{Latex[length=2.5mm]}, thick, draw=negative, dashed},
  % Data-point styles for plots
  observed-dot/.style = {fill=primary-blue, draw=primary-blue, circle, inner sep=1.2pt},
  counterfactual-dot/.style = {fill=white, draw=neutral, circle, inner sep=1.2pt}
}

% -----------------------------------------------------------------------------
% Convenience macros
% -----------------------------------------------------------------------------
\newcommand{\muted}[1]{\textcolor{neutral}{#1}}
\newcommand{\key}[1]{\textcolor{primary-gold}{\textbf{#1}}}
\newcommand{\good}[1]{\textcolor{positive}{#1}}
\newcommand{\bad}[1]{\textcolor{negative}{#1}}

% Standout frame macro: full-bleed dark background for section transitions.
% Beamer-only (uses \begin{frame}/\setbeamercolor/\usebeamerfont) -- do not
% call from an article-class file (e.g. Notes/<CODE>/*-notes.tex). \muted,
% \key, \good, \bad, and \coursecode above are plain xcolor/gdef and are
% safe to use from either class; verified when Notes/article-class support
% was added.
% Expand: \transitionslide{Your transition title}
\newcommand{\transitionslide}[1]{%
  \begingroup
  \setbeamercolor{background canvas}{bg=primary-blue}
  \begin{frame}[plain]
    \centering
    \vfill
    {\usebeamerfont{title}\color{white}\Large #1\par}
    \vfill
  \end{frame}
  \endgroup
}
, before \begin{document})
% via \coursecode{CS401}. Shown in the footer on every slide so a deck's course
% is identifiable without opening the title page. Empty by default.
% -----------------------------------------------------------------------------
\makeatletter
\newcommand{\coursecode}[1]{\gdef\@coursecodetext{#1}}
\gdef\@coursecodetext{}
\makeatother

% -----------------------------------------------------------------------------
% Beamer theme assignments (only apply under Beamer).
%
% We detect Beamer via \@ifclassloaded{beamer}, which is the documented,
% non-internal way. This must be wrapped in \makeatletter/\makeatother
% because \@ifclassloaded contains an @-catcode character.
% -----------------------------------------------------------------------------
\makeatletter
\@ifclassloaded{beamer}{%
  \usefonttheme{professionalfonts}
  \setbeamercolor{structure}{fg=primary-blue}
  \setbeamercolor{frametitle}{fg=primary-blue}
  \setbeamercolor{title}{fg=primary-blue}
  \setbeamercolor{subtitle}{fg=primary-gold}
  \setbeamercolor{author}{fg=primary-blue}
  \setbeamercolor{institute}{fg=neutral}
  \setbeamercolor{date}{fg=primary-gold}
  \setbeamercolor{itemize item}{fg=primary-blue}
  \setbeamercolor{itemize subitem}{fg=primary-gold}
  \setbeamercolor{alerted text}{fg=primary-gold}
  \setbeamercolor{block title}{fg=white, bg=primary-blue}
  \setbeamercolor{block body}{bg=light-bg}
  % Semantic block variants: block=definition/concept (blue), exampleblock=
  % worked example (green/positive), alertblock=key takeaway or caution (gold).
  % Keep to <=2 colored boxes per slide across ALL three variants (INV-7).
  \setbeamercolor{block title example}{fg=white, bg=positive}
  \setbeamercolor{block body example}{bg=light-bg}
  \setbeamercolor{block title alerted}{fg=white, bg=primary-gold}
  \setbeamercolor{block body alerted}{bg=light-bg}

  % ---------------------------------------------------------------------------
  % Formal title page: serif (Libertine, via \rmfamily) for title-page
  % elements only. Body text, itemize, and frametitle are intentionally left
  % on Beamer's sans-serif default -- \usebeamerfont{title} is also what
  % \transitionslide (below) uses for its big centered text, so section
  % transitions pick up the same serif treatment for free.
  % ---------------------------------------------------------------------------
  \setbeamerfont{title}{family=\rmfamily, series=\bfseries, size=\Large}
  \setbeamerfont{subtitle}{family=\rmfamily, shape=\itshape, size=\normalsize}
  \setbeamerfont{author}{family=\rmfamily, size=\normalsize}
  \setbeamerfont{institute}{family=\rmfamily, size=\scriptsize}
  \setbeamerfont{date}{family=\rmfamily, size=\scriptsize}

  \setbeamertemplate{title page}{
    \vfill
    \begin{center}
      {\usebeamerfont{title}\usebeamercolor[fg]{title}\inserttitle\par}
      \vspace{0.15cm}
      {\usebeamerfont{subtitle}\usebeamercolor[fg]{subtitle}\insertsubtitle\par}
      \vspace{0.45cm}
      {\color{primary-gold}\rule{0.42\textwidth}{0.6pt}\par}
      \vspace{0.45cm}
      {\usebeamerfont{author}\usebeamercolor[fg]{author}\insertauthor\par}
      \vspace{0.35cm}
      {\usebeamerfont{institute}\usebeamercolor[fg]{institute}\insertinstitute\par}
      \vspace{0.35cm}
      {\usebeamerfont{date}\usebeamercolor[fg]{date}\insertdate\par}
    \end{center}
    \vfill
  }

  % Minimal footer: course code (left) + page X / N (right), no navigation clutter
  \setbeamertemplate{navigation symbols}{}
  \setbeamertemplate{footline}{%
    \color{neutral}\footnotesize\hspace{0.7em}\@coursecodetext\hfill
    \insertframenumber/\inserttotalframenumber\hspace{0.7em}\vspace{0.3em}}
}{}
\makeatother

% -----------------------------------------------------------------------------
% TikZ setup — libraries the snippet gallery uses
% -----------------------------------------------------------------------------
\usepackage{tikz}
\usetikzlibrary{arrows.meta, positioning, calc, decorations.pathreplacing,
                fit, shapes.geometric, backgrounds}

% Shared TikZ styles — snippets and hand-written diagrams can pick these up
% via `\begin{tikzpicture}[every node/.style={...}]` or by naming specific
% styles (e.g., `\node[dag-node] (X) at (0,0) {X};`).
\tikzset{
  % Node styles for causal DAGs and similar
  dag-node/.style = {
    draw=primary-blue, fill=light-bg, rounded corners=2pt,
    minimum width=1.2cm, minimum height=0.9cm, align=center, font=\small
  },
  decision-node/.style = {
    draw=primary-gold, fill=white, diamond, aspect=1.8,
    minimum width=1.8cm, minimum height=1.2cm, align=center, font=\small,
    inner sep=1pt
  },
  % Edge styles — observed vs counterfactual
  observed-edge/.style = {-{Latex[length=2.5mm]}, thick, draw=jet},
  counterfactual-edge/.style = {-{Latex[length=2.5mm]}, thick, draw=neutral, dashed},
  confound-edge/.style = {-{Latex[length=2.5mm]}, thick, draw=negative, dashed},
  % Data-point styles for plots
  observed-dot/.style = {fill=primary-blue, draw=primary-blue, circle, inner sep=1.2pt},
  counterfactual-dot/.style = {fill=white, draw=neutral, circle, inner sep=1.2pt}
}

% -----------------------------------------------------------------------------
% Convenience macros
% -----------------------------------------------------------------------------
\newcommand{\muted}[1]{\textcolor{neutral}{#1}}
\newcommand{\key}[1]{\textcolor{primary-gold}{\textbf{#1}}}
\newcommand{\good}[1]{\textcolor{positive}{#1}}
\newcommand{\bad}[1]{\textcolor{negative}{#1}}

% Standout frame macro: full-bleed dark background for section transitions.
% Beamer-only (uses \begin{frame}/\setbeamercolor/\usebeamerfont) -- do not
% call from an article-class file (e.g. Notes/<CODE>/*-notes.tex). \muted,
% \key, \good, \bad, and \coursecode above are plain xcolor/gdef and are
% safe to use from either class; verified when Notes/article-class support
% was added.
% Expand: \transitionslide{Your transition title}
\newcommand{\transitionslide}[1]{%
  \begingroup
  \setbeamercolor{background canvas}{bg=primary-blue}
  \begin{frame}[plain]
    \centering
    \vfill
    {\usebeamerfont{title}\color{white}\Large #1\par}
    \vfill
  \end{frame}
  \endgroup
}

% and compiling with TEXINPUTS=../Preambles:$TEXINPUTS (the /compile-latex
% skill does this automatically).
%
% PALETTE CONTRACT. Color names defined here MUST match the names in
% Quarto/theme-template.scss so Beamer and Quarto renderings use the same
% palette. The scripts/check-palette-sync.sh script enforces this.
%
% To customize: edit the HEX values below AND the matching $variable in
% Quarto/theme-template.scss, then run scripts/check-palette-sync.sh.
%
% TITLE PAGE. A formal, serif (Libertine) title-page template is defined in
% the Beamer-only block below (\setbeamertemplate{title page}) -- title,
% subtitle, a thin gold rule, then \author/\institute/\date. Frametitle and
% body text remain on Beamer's sans-serif default. No SCSS counterpart is
% needed for this (font/layout only, no new colors).
% =============================================================================

% -----------------------------------------------------------------------------
% Engine detection + encoding
%
% Our /compile-latex skill uses XeLaTeX, where inputenc is unnecessary and
% can produce encoding warnings. Only load inputenc under pdfTeX.
% -----------------------------------------------------------------------------
\usepackage{iftex}
\ifPDFTeX
  \usepackage[utf8]{inputenc}
\fi

\usepackage{amsmath, amssymb, amsthm}
\usepackage{booktabs}
\usepackage{graphicx}

% -----------------------------------------------------------------------------
% Serif face for the title page (formal/academic look). Linux Libertine is
% XeLaTeX- and pdfTeX-safe and redefines \rmdefault, so \rmfamily picks it up
% everywhere \rmfamily is used explicitly. Body text and frametitles stay on
% Beamer's own sans-serif default (unaffected) -- only the title-page
% \setbeamerfont{...} overrides below opt in to \rmfamily.
% -----------------------------------------------------------------------------
\usepackage{libertine}

% -----------------------------------------------------------------------------
% PALETTE — mirrors Quarto/theme-template.scss variable names.
% Load xcolor and define colors BEFORE anything that references them
% (notably hyperref, hypersetup, and the Beamer color assignments).
% -----------------------------------------------------------------------------
\usepackage{xcolor}

\definecolor{primary-blue}{HTML}{012169}      % headings, accents
\definecolor{primary-gold}{HTML}{B9975B}      % emphasis, borders
\definecolor{highlight-yellow}{HTML}{F2A900}  % markers, alerts
\definecolor{light-bg}{HTML}{E8EDF5}          % title / section backgrounds
\definecolor{jet}{HTML}{1A1A1A}               % body text

% Semantic colors (used in diagrams and annotations)
\definecolor{positive}{HTML}{15803D}          % good / observed / identified
\definecolor{negative}{HTML}{B91C1C}          % bad / problematic / bias
\definecolor{neutral}{HTML}{525252}           % reference / context

% Highlight accents (match .hi-* classes in the SCSS)
\definecolor{hi-slate}{HTML}{314F4F}
\definecolor{hi-green}{HTML}{2E8B57}
\definecolor{hi-red}{HTML}{C41E3A}

% Semantic aliases so skill templates and snippets can write
% \textcolor{accent}{...} without hard-coding "primary-blue".
\colorlet{accent}{primary-blue}
\colorlet{accent2}{primary-gold}

% -----------------------------------------------------------------------------
% Hyperref — loaded AFTER xcolor + palette so link colors resolve.
% -----------------------------------------------------------------------------
\usepackage{hyperref}
\hypersetup{colorlinks=true, linkcolor=primary-blue, urlcolor=primary-blue,
            citecolor=primary-blue, pdfborder={0 0 0}}

% -----------------------------------------------------------------------------
% Course code — set once per deck (after % =============================================================================
% Preambles/header.tex — shared LaTeX/Beamer preamble for this project
%
% Use from a Beamer lecture by prepending:
%     % =============================================================================
% Preambles/header.tex — shared LaTeX/Beamer preamble for this project
%
% Use from a Beamer lecture by prepending:
%     \input{header}
% and compiling with TEXINPUTS=../Preambles:$TEXINPUTS (the /compile-latex
% skill does this automatically).
%
% PALETTE CONTRACT. Color names defined here MUST match the names in
% Quarto/theme-template.scss so Beamer and Quarto renderings use the same
% palette. The scripts/check-palette-sync.sh script enforces this.
%
% To customize: edit the HEX values below AND the matching $variable in
% Quarto/theme-template.scss, then run scripts/check-palette-sync.sh.
%
% TITLE PAGE. A formal, serif (Libertine) title-page template is defined in
% the Beamer-only block below (\setbeamertemplate{title page}) -- title,
% subtitle, a thin gold rule, then \author/\institute/\date. Frametitle and
% body text remain on Beamer's sans-serif default. No SCSS counterpart is
% needed for this (font/layout only, no new colors).
% =============================================================================

% -----------------------------------------------------------------------------
% Engine detection + encoding
%
% Our /compile-latex skill uses XeLaTeX, where inputenc is unnecessary and
% can produce encoding warnings. Only load inputenc under pdfTeX.
% -----------------------------------------------------------------------------
\usepackage{iftex}
\ifPDFTeX
  \usepackage[utf8]{inputenc}
\fi

\usepackage{amsmath, amssymb, amsthm}
\usepackage{booktabs}
\usepackage{graphicx}

% -----------------------------------------------------------------------------
% Serif face for the title page (formal/academic look). Linux Libertine is
% XeLaTeX- and pdfTeX-safe and redefines \rmdefault, so \rmfamily picks it up
% everywhere \rmfamily is used explicitly. Body text and frametitles stay on
% Beamer's own sans-serif default (unaffected) -- only the title-page
% \setbeamerfont{...} overrides below opt in to \rmfamily.
% -----------------------------------------------------------------------------
\usepackage{libertine}

% -----------------------------------------------------------------------------
% PALETTE — mirrors Quarto/theme-template.scss variable names.
% Load xcolor and define colors BEFORE anything that references them
% (notably hyperref, hypersetup, and the Beamer color assignments).
% -----------------------------------------------------------------------------
\usepackage{xcolor}

\definecolor{primary-blue}{HTML}{012169}      % headings, accents
\definecolor{primary-gold}{HTML}{B9975B}      % emphasis, borders
\definecolor{highlight-yellow}{HTML}{F2A900}  % markers, alerts
\definecolor{light-bg}{HTML}{E8EDF5}          % title / section backgrounds
\definecolor{jet}{HTML}{1A1A1A}               % body text

% Semantic colors (used in diagrams and annotations)
\definecolor{positive}{HTML}{15803D}          % good / observed / identified
\definecolor{negative}{HTML}{B91C1C}          % bad / problematic / bias
\definecolor{neutral}{HTML}{525252}           % reference / context

% Highlight accents (match .hi-* classes in the SCSS)
\definecolor{hi-slate}{HTML}{314F4F}
\definecolor{hi-green}{HTML}{2E8B57}
\definecolor{hi-red}{HTML}{C41E3A}

% Semantic aliases so skill templates and snippets can write
% \textcolor{accent}{...} without hard-coding "primary-blue".
\colorlet{accent}{primary-blue}
\colorlet{accent2}{primary-gold}

% -----------------------------------------------------------------------------
% Hyperref — loaded AFTER xcolor + palette so link colors resolve.
% -----------------------------------------------------------------------------
\usepackage{hyperref}
\hypersetup{colorlinks=true, linkcolor=primary-blue, urlcolor=primary-blue,
            citecolor=primary-blue, pdfborder={0 0 0}}

% -----------------------------------------------------------------------------
% Course code — set once per deck (after \input{header}, before \begin{document})
% via \coursecode{CS401}. Shown in the footer on every slide so a deck's course
% is identifiable without opening the title page. Empty by default.
% -----------------------------------------------------------------------------
\makeatletter
\newcommand{\coursecode}[1]{\gdef\@coursecodetext{#1}}
\gdef\@coursecodetext{}
\makeatother

% -----------------------------------------------------------------------------
% Beamer theme assignments (only apply under Beamer).
%
% We detect Beamer via \@ifclassloaded{beamer}, which is the documented,
% non-internal way. This must be wrapped in \makeatletter/\makeatother
% because \@ifclassloaded contains an @-catcode character.
% -----------------------------------------------------------------------------
\makeatletter
\@ifclassloaded{beamer}{%
  \usefonttheme{professionalfonts}
  \setbeamercolor{structure}{fg=primary-blue}
  \setbeamercolor{frametitle}{fg=primary-blue}
  \setbeamercolor{title}{fg=primary-blue}
  \setbeamercolor{subtitle}{fg=primary-gold}
  \setbeamercolor{author}{fg=primary-blue}
  \setbeamercolor{institute}{fg=neutral}
  \setbeamercolor{date}{fg=primary-gold}
  \setbeamercolor{itemize item}{fg=primary-blue}
  \setbeamercolor{itemize subitem}{fg=primary-gold}
  \setbeamercolor{alerted text}{fg=primary-gold}
  \setbeamercolor{block title}{fg=white, bg=primary-blue}
  \setbeamercolor{block body}{bg=light-bg}
  % Semantic block variants: block=definition/concept (blue), exampleblock=
  % worked example (green/positive), alertblock=key takeaway or caution (gold).
  % Keep to <=2 colored boxes per slide across ALL three variants (INV-7).
  \setbeamercolor{block title example}{fg=white, bg=positive}
  \setbeamercolor{block body example}{bg=light-bg}
  \setbeamercolor{block title alerted}{fg=white, bg=primary-gold}
  \setbeamercolor{block body alerted}{bg=light-bg}

  % ---------------------------------------------------------------------------
  % Formal title page: serif (Libertine, via \rmfamily) for title-page
  % elements only. Body text, itemize, and frametitle are intentionally left
  % on Beamer's sans-serif default -- \usebeamerfont{title} is also what
  % \transitionslide (below) uses for its big centered text, so section
  % transitions pick up the same serif treatment for free.
  % ---------------------------------------------------------------------------
  \setbeamerfont{title}{family=\rmfamily, series=\bfseries, size=\Large}
  \setbeamerfont{subtitle}{family=\rmfamily, shape=\itshape, size=\normalsize}
  \setbeamerfont{author}{family=\rmfamily, size=\normalsize}
  \setbeamerfont{institute}{family=\rmfamily, size=\scriptsize}
  \setbeamerfont{date}{family=\rmfamily, size=\scriptsize}

  \setbeamertemplate{title page}{
    \vfill
    \begin{center}
      {\usebeamerfont{title}\usebeamercolor[fg]{title}\inserttitle\par}
      \vspace{0.15cm}
      {\usebeamerfont{subtitle}\usebeamercolor[fg]{subtitle}\insertsubtitle\par}
      \vspace{0.45cm}
      {\color{primary-gold}\rule{0.42\textwidth}{0.6pt}\par}
      \vspace{0.45cm}
      {\usebeamerfont{author}\usebeamercolor[fg]{author}\insertauthor\par}
      \vspace{0.35cm}
      {\usebeamerfont{institute}\usebeamercolor[fg]{institute}\insertinstitute\par}
      \vspace{0.35cm}
      {\usebeamerfont{date}\usebeamercolor[fg]{date}\insertdate\par}
    \end{center}
    \vfill
  }

  % Minimal footer: course code (left) + page X / N (right), no navigation clutter
  \setbeamertemplate{navigation symbols}{}
  \setbeamertemplate{footline}{%
    \color{neutral}\footnotesize\hspace{0.7em}\@coursecodetext\hfill
    \insertframenumber/\inserttotalframenumber\hspace{0.7em}\vspace{0.3em}}
}{}
\makeatother

% -----------------------------------------------------------------------------
% TikZ setup — libraries the snippet gallery uses
% -----------------------------------------------------------------------------
\usepackage{tikz}
\usetikzlibrary{arrows.meta, positioning, calc, decorations.pathreplacing,
                fit, shapes.geometric, backgrounds}

% Shared TikZ styles — snippets and hand-written diagrams can pick these up
% via `\begin{tikzpicture}[every node/.style={...}]` or by naming specific
% styles (e.g., `\node[dag-node] (X) at (0,0) {X};`).
\tikzset{
  % Node styles for causal DAGs and similar
  dag-node/.style = {
    draw=primary-blue, fill=light-bg, rounded corners=2pt,
    minimum width=1.2cm, minimum height=0.9cm, align=center, font=\small
  },
  decision-node/.style = {
    draw=primary-gold, fill=white, diamond, aspect=1.8,
    minimum width=1.8cm, minimum height=1.2cm, align=center, font=\small,
    inner sep=1pt
  },
  % Edge styles — observed vs counterfactual
  observed-edge/.style = {-{Latex[length=2.5mm]}, thick, draw=jet},
  counterfactual-edge/.style = {-{Latex[length=2.5mm]}, thick, draw=neutral, dashed},
  confound-edge/.style = {-{Latex[length=2.5mm]}, thick, draw=negative, dashed},
  % Data-point styles for plots
  observed-dot/.style = {fill=primary-blue, draw=primary-blue, circle, inner sep=1.2pt},
  counterfactual-dot/.style = {fill=white, draw=neutral, circle, inner sep=1.2pt}
}

% -----------------------------------------------------------------------------
% Convenience macros
% -----------------------------------------------------------------------------
\newcommand{\muted}[1]{\textcolor{neutral}{#1}}
\newcommand{\key}[1]{\textcolor{primary-gold}{\textbf{#1}}}
\newcommand{\good}[1]{\textcolor{positive}{#1}}
\newcommand{\bad}[1]{\textcolor{negative}{#1}}

% Standout frame macro: full-bleed dark background for section transitions.
% Beamer-only (uses \begin{frame}/\setbeamercolor/\usebeamerfont) -- do not
% call from an article-class file (e.g. Notes/<CODE>/*-notes.tex). \muted,
% \key, \good, \bad, and \coursecode above are plain xcolor/gdef and are
% safe to use from either class; verified when Notes/article-class support
% was added.
% Expand: \transitionslide{Your transition title}
\newcommand{\transitionslide}[1]{%
  \begingroup
  \setbeamercolor{background canvas}{bg=primary-blue}
  \begin{frame}[plain]
    \centering
    \vfill
    {\usebeamerfont{title}\color{white}\Large #1\par}
    \vfill
  \end{frame}
  \endgroup
}

% and compiling with TEXINPUTS=../Preambles:$TEXINPUTS (the /compile-latex
% skill does this automatically).
%
% PALETTE CONTRACT. Color names defined here MUST match the names in
% Quarto/theme-template.scss so Beamer and Quarto renderings use the same
% palette. The scripts/check-palette-sync.sh script enforces this.
%
% To customize: edit the HEX values below AND the matching $variable in
% Quarto/theme-template.scss, then run scripts/check-palette-sync.sh.
%
% TITLE PAGE. A formal, serif (Libertine) title-page template is defined in
% the Beamer-only block below (\setbeamertemplate{title page}) -- title,
% subtitle, a thin gold rule, then \author/\institute/\date. Frametitle and
% body text remain on Beamer's sans-serif default. No SCSS counterpart is
% needed for this (font/layout only, no new colors).
% =============================================================================

% -----------------------------------------------------------------------------
% Engine detection + encoding
%
% Our /compile-latex skill uses XeLaTeX, where inputenc is unnecessary and
% can produce encoding warnings. Only load inputenc under pdfTeX.
% -----------------------------------------------------------------------------
\usepackage{iftex}
\ifPDFTeX
  \usepackage[utf8]{inputenc}
\fi

\usepackage{amsmath, amssymb, amsthm}
\usepackage{booktabs}
\usepackage{graphicx}

% -----------------------------------------------------------------------------
% Serif face for the title page (formal/academic look). Linux Libertine is
% XeLaTeX- and pdfTeX-safe and redefines \rmdefault, so \rmfamily picks it up
% everywhere \rmfamily is used explicitly. Body text and frametitles stay on
% Beamer's own sans-serif default (unaffected) -- only the title-page
% \setbeamerfont{...} overrides below opt in to \rmfamily.
% -----------------------------------------------------------------------------
\usepackage{libertine}

% -----------------------------------------------------------------------------
% PALETTE — mirrors Quarto/theme-template.scss variable names.
% Load xcolor and define colors BEFORE anything that references them
% (notably hyperref, hypersetup, and the Beamer color assignments).
% -----------------------------------------------------------------------------
\usepackage{xcolor}

\definecolor{primary-blue}{HTML}{012169}      % headings, accents
\definecolor{primary-gold}{HTML}{B9975B}      % emphasis, borders
\definecolor{highlight-yellow}{HTML}{F2A900}  % markers, alerts
\definecolor{light-bg}{HTML}{E8EDF5}          % title / section backgrounds
\definecolor{jet}{HTML}{1A1A1A}               % body text

% Semantic colors (used in diagrams and annotations)
\definecolor{positive}{HTML}{15803D}          % good / observed / identified
\definecolor{negative}{HTML}{B91C1C}          % bad / problematic / bias
\definecolor{neutral}{HTML}{525252}           % reference / context

% Highlight accents (match .hi-* classes in the SCSS)
\definecolor{hi-slate}{HTML}{314F4F}
\definecolor{hi-green}{HTML}{2E8B57}
\definecolor{hi-red}{HTML}{C41E3A}

% Semantic aliases so skill templates and snippets can write
% \textcolor{accent}{...} without hard-coding "primary-blue".
\colorlet{accent}{primary-blue}
\colorlet{accent2}{primary-gold}

% -----------------------------------------------------------------------------
% Hyperref — loaded AFTER xcolor + palette so link colors resolve.
% -----------------------------------------------------------------------------
\usepackage{hyperref}
\hypersetup{colorlinks=true, linkcolor=primary-blue, urlcolor=primary-blue,
            citecolor=primary-blue, pdfborder={0 0 0}}

% -----------------------------------------------------------------------------
% Course code — set once per deck (after % =============================================================================
% Preambles/header.tex — shared LaTeX/Beamer preamble for this project
%
% Use from a Beamer lecture by prepending:
%     \input{header}
% and compiling with TEXINPUTS=../Preambles:$TEXINPUTS (the /compile-latex
% skill does this automatically).
%
% PALETTE CONTRACT. Color names defined here MUST match the names in
% Quarto/theme-template.scss so Beamer and Quarto renderings use the same
% palette. The scripts/check-palette-sync.sh script enforces this.
%
% To customize: edit the HEX values below AND the matching $variable in
% Quarto/theme-template.scss, then run scripts/check-palette-sync.sh.
%
% TITLE PAGE. A formal, serif (Libertine) title-page template is defined in
% the Beamer-only block below (\setbeamertemplate{title page}) -- title,
% subtitle, a thin gold rule, then \author/\institute/\date. Frametitle and
% body text remain on Beamer's sans-serif default. No SCSS counterpart is
% needed for this (font/layout only, no new colors).
% =============================================================================

% -----------------------------------------------------------------------------
% Engine detection + encoding
%
% Our /compile-latex skill uses XeLaTeX, where inputenc is unnecessary and
% can produce encoding warnings. Only load inputenc under pdfTeX.
% -----------------------------------------------------------------------------
\usepackage{iftex}
\ifPDFTeX
  \usepackage[utf8]{inputenc}
\fi

\usepackage{amsmath, amssymb, amsthm}
\usepackage{booktabs}
\usepackage{graphicx}

% -----------------------------------------------------------------------------
% Serif face for the title page (formal/academic look). Linux Libertine is
% XeLaTeX- and pdfTeX-safe and redefines \rmdefault, so \rmfamily picks it up
% everywhere \rmfamily is used explicitly. Body text and frametitles stay on
% Beamer's own sans-serif default (unaffected) -- only the title-page
% \setbeamerfont{...} overrides below opt in to \rmfamily.
% -----------------------------------------------------------------------------
\usepackage{libertine}

% -----------------------------------------------------------------------------
% PALETTE — mirrors Quarto/theme-template.scss variable names.
% Load xcolor and define colors BEFORE anything that references them
% (notably hyperref, hypersetup, and the Beamer color assignments).
% -----------------------------------------------------------------------------
\usepackage{xcolor}

\definecolor{primary-blue}{HTML}{012169}      % headings, accents
\definecolor{primary-gold}{HTML}{B9975B}      % emphasis, borders
\definecolor{highlight-yellow}{HTML}{F2A900}  % markers, alerts
\definecolor{light-bg}{HTML}{E8EDF5}          % title / section backgrounds
\definecolor{jet}{HTML}{1A1A1A}               % body text

% Semantic colors (used in diagrams and annotations)
\definecolor{positive}{HTML}{15803D}          % good / observed / identified
\definecolor{negative}{HTML}{B91C1C}          % bad / problematic / bias
\definecolor{neutral}{HTML}{525252}           % reference / context

% Highlight accents (match .hi-* classes in the SCSS)
\definecolor{hi-slate}{HTML}{314F4F}
\definecolor{hi-green}{HTML}{2E8B57}
\definecolor{hi-red}{HTML}{C41E3A}

% Semantic aliases so skill templates and snippets can write
% \textcolor{accent}{...} without hard-coding "primary-blue".
\colorlet{accent}{primary-blue}
\colorlet{accent2}{primary-gold}

% -----------------------------------------------------------------------------
% Hyperref — loaded AFTER xcolor + palette so link colors resolve.
% -----------------------------------------------------------------------------
\usepackage{hyperref}
\hypersetup{colorlinks=true, linkcolor=primary-blue, urlcolor=primary-blue,
            citecolor=primary-blue, pdfborder={0 0 0}}

% -----------------------------------------------------------------------------
% Course code — set once per deck (after \input{header}, before \begin{document})
% via \coursecode{CS401}. Shown in the footer on every slide so a deck's course
% is identifiable without opening the title page. Empty by default.
% -----------------------------------------------------------------------------
\makeatletter
\newcommand{\coursecode}[1]{\gdef\@coursecodetext{#1}}
\gdef\@coursecodetext{}
\makeatother

% -----------------------------------------------------------------------------
% Beamer theme assignments (only apply under Beamer).
%
% We detect Beamer via \@ifclassloaded{beamer}, which is the documented,
% non-internal way. This must be wrapped in \makeatletter/\makeatother
% because \@ifclassloaded contains an @-catcode character.
% -----------------------------------------------------------------------------
\makeatletter
\@ifclassloaded{beamer}{%
  \usefonttheme{professionalfonts}
  \setbeamercolor{structure}{fg=primary-blue}
  \setbeamercolor{frametitle}{fg=primary-blue}
  \setbeamercolor{title}{fg=primary-blue}
  \setbeamercolor{subtitle}{fg=primary-gold}
  \setbeamercolor{author}{fg=primary-blue}
  \setbeamercolor{institute}{fg=neutral}
  \setbeamercolor{date}{fg=primary-gold}
  \setbeamercolor{itemize item}{fg=primary-blue}
  \setbeamercolor{itemize subitem}{fg=primary-gold}
  \setbeamercolor{alerted text}{fg=primary-gold}
  \setbeamercolor{block title}{fg=white, bg=primary-blue}
  \setbeamercolor{block body}{bg=light-bg}
  % Semantic block variants: block=definition/concept (blue), exampleblock=
  % worked example (green/positive), alertblock=key takeaway or caution (gold).
  % Keep to <=2 colored boxes per slide across ALL three variants (INV-7).
  \setbeamercolor{block title example}{fg=white, bg=positive}
  \setbeamercolor{block body example}{bg=light-bg}
  \setbeamercolor{block title alerted}{fg=white, bg=primary-gold}
  \setbeamercolor{block body alerted}{bg=light-bg}

  % ---------------------------------------------------------------------------
  % Formal title page: serif (Libertine, via \rmfamily) for title-page
  % elements only. Body text, itemize, and frametitle are intentionally left
  % on Beamer's sans-serif default -- \usebeamerfont{title} is also what
  % \transitionslide (below) uses for its big centered text, so section
  % transitions pick up the same serif treatment for free.
  % ---------------------------------------------------------------------------
  \setbeamerfont{title}{family=\rmfamily, series=\bfseries, size=\Large}
  \setbeamerfont{subtitle}{family=\rmfamily, shape=\itshape, size=\normalsize}
  \setbeamerfont{author}{family=\rmfamily, size=\normalsize}
  \setbeamerfont{institute}{family=\rmfamily, size=\scriptsize}
  \setbeamerfont{date}{family=\rmfamily, size=\scriptsize}

  \setbeamertemplate{title page}{
    \vfill
    \begin{center}
      {\usebeamerfont{title}\usebeamercolor[fg]{title}\inserttitle\par}
      \vspace{0.15cm}
      {\usebeamerfont{subtitle}\usebeamercolor[fg]{subtitle}\insertsubtitle\par}
      \vspace{0.45cm}
      {\color{primary-gold}\rule{0.42\textwidth}{0.6pt}\par}
      \vspace{0.45cm}
      {\usebeamerfont{author}\usebeamercolor[fg]{author}\insertauthor\par}
      \vspace{0.35cm}
      {\usebeamerfont{institute}\usebeamercolor[fg]{institute}\insertinstitute\par}
      \vspace{0.35cm}
      {\usebeamerfont{date}\usebeamercolor[fg]{date}\insertdate\par}
    \end{center}
    \vfill
  }

  % Minimal footer: course code (left) + page X / N (right), no navigation clutter
  \setbeamertemplate{navigation symbols}{}
  \setbeamertemplate{footline}{%
    \color{neutral}\footnotesize\hspace{0.7em}\@coursecodetext\hfill
    \insertframenumber/\inserttotalframenumber\hspace{0.7em}\vspace{0.3em}}
}{}
\makeatother

% -----------------------------------------------------------------------------
% TikZ setup — libraries the snippet gallery uses
% -----------------------------------------------------------------------------
\usepackage{tikz}
\usetikzlibrary{arrows.meta, positioning, calc, decorations.pathreplacing,
                fit, shapes.geometric, backgrounds}

% Shared TikZ styles — snippets and hand-written diagrams can pick these up
% via `\begin{tikzpicture}[every node/.style={...}]` or by naming specific
% styles (e.g., `\node[dag-node] (X) at (0,0) {X};`).
\tikzset{
  % Node styles for causal DAGs and similar
  dag-node/.style = {
    draw=primary-blue, fill=light-bg, rounded corners=2pt,
    minimum width=1.2cm, minimum height=0.9cm, align=center, font=\small
  },
  decision-node/.style = {
    draw=primary-gold, fill=white, diamond, aspect=1.8,
    minimum width=1.8cm, minimum height=1.2cm, align=center, font=\small,
    inner sep=1pt
  },
  % Edge styles — observed vs counterfactual
  observed-edge/.style = {-{Latex[length=2.5mm]}, thick, draw=jet},
  counterfactual-edge/.style = {-{Latex[length=2.5mm]}, thick, draw=neutral, dashed},
  confound-edge/.style = {-{Latex[length=2.5mm]}, thick, draw=negative, dashed},
  % Data-point styles for plots
  observed-dot/.style = {fill=primary-blue, draw=primary-blue, circle, inner sep=1.2pt},
  counterfactual-dot/.style = {fill=white, draw=neutral, circle, inner sep=1.2pt}
}

% -----------------------------------------------------------------------------
% Convenience macros
% -----------------------------------------------------------------------------
\newcommand{\muted}[1]{\textcolor{neutral}{#1}}
\newcommand{\key}[1]{\textcolor{primary-gold}{\textbf{#1}}}
\newcommand{\good}[1]{\textcolor{positive}{#1}}
\newcommand{\bad}[1]{\textcolor{negative}{#1}}

% Standout frame macro: full-bleed dark background for section transitions.
% Beamer-only (uses \begin{frame}/\setbeamercolor/\usebeamerfont) -- do not
% call from an article-class file (e.g. Notes/<CODE>/*-notes.tex). \muted,
% \key, \good, \bad, and \coursecode above are plain xcolor/gdef and are
% safe to use from either class; verified when Notes/article-class support
% was added.
% Expand: \transitionslide{Your transition title}
\newcommand{\transitionslide}[1]{%
  \begingroup
  \setbeamercolor{background canvas}{bg=primary-blue}
  \begin{frame}[plain]
    \centering
    \vfill
    {\usebeamerfont{title}\color{white}\Large #1\par}
    \vfill
  \end{frame}
  \endgroup
}
, before \begin{document})
% via \coursecode{CS401}. Shown in the footer on every slide so a deck's course
% is identifiable without opening the title page. Empty by default.
% -----------------------------------------------------------------------------
\makeatletter
\newcommand{\coursecode}[1]{\gdef\@coursecodetext{#1}}
\gdef\@coursecodetext{}
\makeatother

% -----------------------------------------------------------------------------
% Beamer theme assignments (only apply under Beamer).
%
% We detect Beamer via \@ifclassloaded{beamer}, which is the documented,
% non-internal way. This must be wrapped in \makeatletter/\makeatother
% because \@ifclassloaded contains an @-catcode character.
% -----------------------------------------------------------------------------
\makeatletter
\@ifclassloaded{beamer}{%
  \usefonttheme{professionalfonts}
  \setbeamercolor{structure}{fg=primary-blue}
  \setbeamercolor{frametitle}{fg=primary-blue}
  \setbeamercolor{title}{fg=primary-blue}
  \setbeamercolor{subtitle}{fg=primary-gold}
  \setbeamercolor{author}{fg=primary-blue}
  \setbeamercolor{institute}{fg=neutral}
  \setbeamercolor{date}{fg=primary-gold}
  \setbeamercolor{itemize item}{fg=primary-blue}
  \setbeamercolor{itemize subitem}{fg=primary-gold}
  \setbeamercolor{alerted text}{fg=primary-gold}
  \setbeamercolor{block title}{fg=white, bg=primary-blue}
  \setbeamercolor{block body}{bg=light-bg}
  % Semantic block variants: block=definition/concept (blue), exampleblock=
  % worked example (green/positive), alertblock=key takeaway or caution (gold).
  % Keep to <=2 colored boxes per slide across ALL three variants (INV-7).
  \setbeamercolor{block title example}{fg=white, bg=positive}
  \setbeamercolor{block body example}{bg=light-bg}
  \setbeamercolor{block title alerted}{fg=white, bg=primary-gold}
  \setbeamercolor{block body alerted}{bg=light-bg}

  % ---------------------------------------------------------------------------
  % Formal title page: serif (Libertine, via \rmfamily) for title-page
  % elements only. Body text, itemize, and frametitle are intentionally left
  % on Beamer's sans-serif default -- \usebeamerfont{title} is also what
  % \transitionslide (below) uses for its big centered text, so section
  % transitions pick up the same serif treatment for free.
  % ---------------------------------------------------------------------------
  \setbeamerfont{title}{family=\rmfamily, series=\bfseries, size=\Large}
  \setbeamerfont{subtitle}{family=\rmfamily, shape=\itshape, size=\normalsize}
  \setbeamerfont{author}{family=\rmfamily, size=\normalsize}
  \setbeamerfont{institute}{family=\rmfamily, size=\scriptsize}
  \setbeamerfont{date}{family=\rmfamily, size=\scriptsize}

  \setbeamertemplate{title page}{
    \vfill
    \begin{center}
      {\usebeamerfont{title}\usebeamercolor[fg]{title}\inserttitle\par}
      \vspace{0.15cm}
      {\usebeamerfont{subtitle}\usebeamercolor[fg]{subtitle}\insertsubtitle\par}
      \vspace{0.45cm}
      {\color{primary-gold}\rule{0.42\textwidth}{0.6pt}\par}
      \vspace{0.45cm}
      {\usebeamerfont{author}\usebeamercolor[fg]{author}\insertauthor\par}
      \vspace{0.35cm}
      {\usebeamerfont{institute}\usebeamercolor[fg]{institute}\insertinstitute\par}
      \vspace{0.35cm}
      {\usebeamerfont{date}\usebeamercolor[fg]{date}\insertdate\par}
    \end{center}
    \vfill
  }

  % Minimal footer: course code (left) + page X / N (right), no navigation clutter
  \setbeamertemplate{navigation symbols}{}
  \setbeamertemplate{footline}{%
    \color{neutral}\footnotesize\hspace{0.7em}\@coursecodetext\hfill
    \insertframenumber/\inserttotalframenumber\hspace{0.7em}\vspace{0.3em}}
}{}
\makeatother

% -----------------------------------------------------------------------------
% TikZ setup — libraries the snippet gallery uses
% -----------------------------------------------------------------------------
\usepackage{tikz}
\usetikzlibrary{arrows.meta, positioning, calc, decorations.pathreplacing,
                fit, shapes.geometric, backgrounds}

% Shared TikZ styles — snippets and hand-written diagrams can pick these up
% via `\begin{tikzpicture}[every node/.style={...}]` or by naming specific
% styles (e.g., `\node[dag-node] (X) at (0,0) {X};`).
\tikzset{
  % Node styles for causal DAGs and similar
  dag-node/.style = {
    draw=primary-blue, fill=light-bg, rounded corners=2pt,
    minimum width=1.2cm, minimum height=0.9cm, align=center, font=\small
  },
  decision-node/.style = {
    draw=primary-gold, fill=white, diamond, aspect=1.8,
    minimum width=1.8cm, minimum height=1.2cm, align=center, font=\small,
    inner sep=1pt
  },
  % Edge styles — observed vs counterfactual
  observed-edge/.style = {-{Latex[length=2.5mm]}, thick, draw=jet},
  counterfactual-edge/.style = {-{Latex[length=2.5mm]}, thick, draw=neutral, dashed},
  confound-edge/.style = {-{Latex[length=2.5mm]}, thick, draw=negative, dashed},
  % Data-point styles for plots
  observed-dot/.style = {fill=primary-blue, draw=primary-blue, circle, inner sep=1.2pt},
  counterfactual-dot/.style = {fill=white, draw=neutral, circle, inner sep=1.2pt}
}

% -----------------------------------------------------------------------------
% Convenience macros
% -----------------------------------------------------------------------------
\newcommand{\muted}[1]{\textcolor{neutral}{#1}}
\newcommand{\key}[1]{\textcolor{primary-gold}{\textbf{#1}}}
\newcommand{\good}[1]{\textcolor{positive}{#1}}
\newcommand{\bad}[1]{\textcolor{negative}{#1}}

% Standout frame macro: full-bleed dark background for section transitions.
% Beamer-only (uses \begin{frame}/\setbeamercolor/\usebeamerfont) -- do not
% call from an article-class file (e.g. Notes/<CODE>/*-notes.tex). \muted,
% \key, \good, \bad, and \coursecode above are plain xcolor/gdef and are
% safe to use from either class; verified when Notes/article-class support
% was added.
% Expand: \transitionslide{Your transition title}
\newcommand{\transitionslide}[1]{%
  \begingroup
  \setbeamercolor{background canvas}{bg=primary-blue}
  \begin{frame}[plain]
    \centering
    \vfill
    {\usebeamerfont{title}\color{white}\Large #1\par}
    \vfill
  \end{frame}
  \endgroup
}
, before \begin{document})
% via \coursecode{CS401}. Shown in the footer on every slide so a deck's course
% is identifiable without opening the title page. Empty by default.
% -----------------------------------------------------------------------------
\makeatletter
\newcommand{\coursecode}[1]{\gdef\@coursecodetext{#1}}
\gdef\@coursecodetext{}
\makeatother

% -----------------------------------------------------------------------------
% Beamer theme assignments (only apply under Beamer).
%
% We detect Beamer via \@ifclassloaded{beamer}, which is the documented,
% non-internal way. This must be wrapped in \makeatletter/\makeatother
% because \@ifclassloaded contains an @-catcode character.
% -----------------------------------------------------------------------------
\makeatletter
\@ifclassloaded{beamer}{%
  \usefonttheme{professionalfonts}
  \setbeamercolor{structure}{fg=primary-blue}
  \setbeamercolor{frametitle}{fg=primary-blue}
  \setbeamercolor{title}{fg=primary-blue}
  \setbeamercolor{subtitle}{fg=primary-gold}
  \setbeamercolor{author}{fg=primary-blue}
  \setbeamercolor{institute}{fg=neutral}
  \setbeamercolor{date}{fg=primary-gold}
  \setbeamercolor{itemize item}{fg=primary-blue}
  \setbeamercolor{itemize subitem}{fg=primary-gold}
  \setbeamercolor{alerted text}{fg=primary-gold}
  \setbeamercolor{block title}{fg=white, bg=primary-blue}
  \setbeamercolor{block body}{bg=light-bg}
  % Semantic block variants: block=definition/concept (blue), exampleblock=
  % worked example (green/positive), alertblock=key takeaway or caution (gold).
  % Keep to <=2 colored boxes per slide across ALL three variants (INV-7).
  \setbeamercolor{block title example}{fg=white, bg=positive}
  \setbeamercolor{block body example}{bg=light-bg}
  \setbeamercolor{block title alerted}{fg=white, bg=primary-gold}
  \setbeamercolor{block body alerted}{bg=light-bg}

  % ---------------------------------------------------------------------------
  % Formal title page: serif (Libertine, via \rmfamily) for title-page
  % elements only. Body text, itemize, and frametitle are intentionally left
  % on Beamer's sans-serif default -- \usebeamerfont{title} is also what
  % \transitionslide (below) uses for its big centered text, so section
  % transitions pick up the same serif treatment for free.
  % ---------------------------------------------------------------------------
  \setbeamerfont{title}{family=\rmfamily, series=\bfseries, size=\Large}
  \setbeamerfont{subtitle}{family=\rmfamily, shape=\itshape, size=\normalsize}
  \setbeamerfont{author}{family=\rmfamily, size=\normalsize}
  \setbeamerfont{institute}{family=\rmfamily, size=\scriptsize}
  \setbeamerfont{date}{family=\rmfamily, size=\scriptsize}

  \setbeamertemplate{title page}{
    \vfill
    \begin{center}
      {\usebeamerfont{title}\usebeamercolor[fg]{title}\inserttitle\par}
      \vspace{0.15cm}
      {\usebeamerfont{subtitle}\usebeamercolor[fg]{subtitle}\insertsubtitle\par}
      \vspace{0.45cm}
      {\color{primary-gold}\rule{0.42\textwidth}{0.6pt}\par}
      \vspace{0.45cm}
      {\usebeamerfont{author}\usebeamercolor[fg]{author}\insertauthor\par}
      \vspace{0.35cm}
      {\usebeamerfont{institute}\usebeamercolor[fg]{institute}\insertinstitute\par}
      \vspace{0.35cm}
      {\usebeamerfont{date}\usebeamercolor[fg]{date}\insertdate\par}
    \end{center}
    \vfill
  }

  % Minimal footer: course code (left) + page X / N (right), no navigation clutter
  \setbeamertemplate{navigation symbols}{}
  \setbeamertemplate{footline}{%
    \color{neutral}\footnotesize\hspace{0.7em}\@coursecodetext\hfill
    \insertframenumber/\inserttotalframenumber\hspace{0.7em}\vspace{0.3em}}
}{}
\makeatother

% -----------------------------------------------------------------------------
% TikZ setup — libraries the snippet gallery uses
% -----------------------------------------------------------------------------
\usepackage{tikz}
\usetikzlibrary{arrows.meta, positioning, calc, decorations.pathreplacing,
                fit, shapes.geometric, backgrounds}

% Shared TikZ styles — snippets and hand-written diagrams can pick these up
% via `\begin{tikzpicture}[every node/.style={...}]` or by naming specific
% styles (e.g., `\node[dag-node] (X) at (0,0) {X};`).
\tikzset{
  % Node styles for causal DAGs and similar
  dag-node/.style = {
    draw=primary-blue, fill=light-bg, rounded corners=2pt,
    minimum width=1.2cm, minimum height=0.9cm, align=center, font=\small
  },
  decision-node/.style = {
    draw=primary-gold, fill=white, diamond, aspect=1.8,
    minimum width=1.8cm, minimum height=1.2cm, align=center, font=\small,
    inner sep=1pt
  },
  % Edge styles — observed vs counterfactual
  observed-edge/.style = {-{Latex[length=2.5mm]}, thick, draw=jet},
  counterfactual-edge/.style = {-{Latex[length=2.5mm]}, thick, draw=neutral, dashed},
  confound-edge/.style = {-{Latex[length=2.5mm]}, thick, draw=negative, dashed},
  % Data-point styles for plots
  observed-dot/.style = {fill=primary-blue, draw=primary-blue, circle, inner sep=1.2pt},
  counterfactual-dot/.style = {fill=white, draw=neutral, circle, inner sep=1.2pt}
}

% -----------------------------------------------------------------------------
% Convenience macros
% -----------------------------------------------------------------------------
\newcommand{\muted}[1]{\textcolor{neutral}{#1}}
\newcommand{\key}[1]{\textcolor{primary-gold}{\textbf{#1}}}
\newcommand{\good}[1]{\textcolor{positive}{#1}}
\newcommand{\bad}[1]{\textcolor{negative}{#1}}

% Standout frame macro: full-bleed dark background for section transitions.
% Beamer-only (uses \begin{frame}/\setbeamercolor/\usebeamerfont) -- do not
% call from an article-class file (e.g. Notes/<CODE>/*-notes.tex). \muted,
% \key, \good, \bad, and \coursecode above are plain xcolor/gdef and are
% safe to use from either class; verified when Notes/article-class support
% was added.
% Expand: \transitionslide{Your transition title}
\newcommand{\transitionslide}[1]{%
  \begingroup
  \setbeamercolor{background canvas}{bg=primary-blue}
  \begin{frame}[plain]
    \centering
    \vfill
    {\usebeamerfont{title}\color{white}\Large #1\par}
    \vfill
  \end{frame}
  \endgroup
}

\coursecode{CS301}
\renewcommand{\thesection}{3.\arabic{section}}

\title{Coding Assignment 3: One Array, One Index, Two Jobs}
\author{Auqib Hamid Lone\\[0.3em]
  \emph{Department of Computer Science \& Engineering}, \emph{GCET Ganderbal}}
\date{\today}

\begin{document}
\maketitle

\noindent\textbf{Due:} two weeks from issue. There is no automatic late penalty --- if you need more time, ask before the deadline and an extension will be granted (see the course policy in the syllabus).

\section{Why this problem}

Week 3 gave the calculator two things: a row of cells that finds \texttt{A[i]} by arithmetic, and a nesting discipline --- \texttt{push}, \texttt{pop}, \texttt{peek} over one array and one index \texttt{top} --- that checks brackets and evaluates postfix in a single left-to-right pass each. The lecture stated the contract but left the \emph{code} for you to write, with one deliberate simplification you must now unlearn: on the slides, \texttt{pop} on an empty stack returns $-1$ to keep the slide readable. That sentinel cannot tell ``the stack held $-1$'' from ``the stack was empty'' --- a real stack must say which happened. This assignment is that real stack, plus the two jobs the lecture promised it could do: Lab P1 (parenthesis check) and Lab P2 (postfix evaluation), both running on your array implementation.

\section{Your task}

Implement the Stack ADT in C as a \textbf{fixed array of \texttt{int} with an explicit status code on every fallible operation}. The contract is given to you in \texttt{stack.h} --- \emph{do not modify it}. Copy \texttt{starter\_stack.c} to \texttt{stack.c} and fill in all eight functions:

\begin{verbatim}
void        stack_init(IntStack *s);
int         stack_is_empty(const IntStack *s);
int         stack_size(const IntStack *s);
StackStatus stack_push(IntStack *s, int x);      /* STACK_OK / STACK_OVERFLOW */
StackStatus stack_pop(IntStack *s, int *out);    /* STACK_OK / STACK_UNDERFLOW */
StackStatus stack_peek(const IntStack *s, int *out);
int         is_balanced(const char *s);          /* P1: 1 if balanced, else 0 */
EvalStatus  eval_postfix(const char *expr, int *out);  /* P2: status + answer */
\end{verbatim}

\textbf{The stack is one array plus \texttt{top}.} \texttt{stack\_init} sets \texttt{top} to $-1$ (empty). \texttt{push} refuses with \texttt{STACK\_OVERFLOW} when \texttt{top + 1 == STACK\_MAX} and changes nothing; \texttt{pop} and \texttt{peek} report \texttt{STACK\_UNDERFLOW} on an empty stack and leave \texttt{*out} unchanged. \texttt{peek} never removes anything. P1 pushes openers as their character codes into the same \texttt{int} cells --- no second stack needed.

\textbf{P1 checks nesting, not content.} Push every \texttt{(}, \texttt{[}, \texttt{\{}; on a closer, return $0$ if the stack is empty or the top is not its match, else pop; every other character (letters, digits, spaces, operators) is ignored, so \texttt{(a+b)*[c]} is checked as \texttt{()[]}. Balanced iff the stack ends empty. The empty string is balanced; \texttt{NULL} returns $0$.

\textbf{P2 evaluates left to right.} Integer literals (optional leading \texttt{-}, no leading \texttt{+}, must fit in \texttt{int}) are pushed; on an operator, \texttt{b} pops \emph{first} and \texttt{a} second, then push $\text{apply}(\text{op}, a, b)$ --- swapping the order breaks \texttt{-}, \texttt{/}, and \texttt{\%}. Tokens are separated by spaces or tabs. On \emph{any} non-\texttt{OK} status, \texttt{*out} is left unchanged. Read the six \texttt{EvalStatus} codes in \texttt{stack.h} before coding --- each error in the table below maps to exactly one of them.

\section{Constraints}

\begin{itemize}
  \item Representation must be the \texttt{IntStack} array in \texttt{stack.h} --- no linked lists, no heap allocation, no global/static mutable state, no \texttt{strtok} shortcuts that hide the token scan.
  \item Error paths use the explicit status codes --- never the lecture's $-1$ sentinel. A submission whose \texttt{pop} returns $-1$-as-error fails the hidden suite's sentinel-trap tests (pushing the \emph{value} $-1$ and popping it back must succeed with \texttt{STACK\_OK}).
  \item Correctness is graded against the contract; your code must compile with \texttt{gcc -Wall -Wextra} without errors or warnings.
  \item Complexity: \texttt{push}, \texttt{pop}, \texttt{peek} are $O(1)$ worst case each; \texttt{is\_balanced} and \texttt{eval\_postfix} each make one left-to-right pass at $O(n)$ time worst case, holding at most \texttt{STACK\_MAX} integers.
\end{itemize}

\section{Worked examples (these are the visible tests)}

\textbf{Example 1 --- the stack itself.} After \texttt{init}: empty, size $0$. \texttt{push(10)}, \texttt{push(20)} both return \texttt{STACK\_OK}; \texttt{peek} gives $20$ and size stays $2$; \texttt{pop} gives $20$, then $10$; \texttt{pop} on the empty stack returns \texttt{STACK\_UNDERFLOW}.

\textbf{Example 2 --- P1.} \texttt{is\_balanced("")} is $1$. \texttt{is\_balanced("(\{[]\})")} is $1$. \texttt{is\_balanced("(\{[\})")} is $0$ (the \texttt{\}} meets \texttt{[}). \texttt{is\_balanced("((())")} is $0$ (an opener never closes).

\textbf{Example 3 --- P2.} \texttt{eval\_postfix("6 2 3 + * 4 -")} returns \texttt{EVAL\_OK} with answer $26$ ($2+3=5$, $6\times5=30$, $30-4=26$). \texttt{eval\_postfix("20 4 / 3 -")} returns \texttt{EVAL\_OK} with answer $2$. \texttt{eval\_postfix("5 0 /")} returns \texttt{EVAL\_DIV\_BY\_ZERO}. \texttt{eval\_postfix("2 +")} returns \texttt{EVAL\_TOO\_FEW\_OPERANDS}.

\section{What to submit}

\begin{enumerate}
  \item \texttt{stack.c} --- your completed implementation (edit \texttt{starter\_stack.c}).
  \item \texttt{stack.h} --- unchanged, but include it so the submission compiles standalone.
\end{enumerate}

\section{Getting started}

\begin{enumerate}
  \item Copy \texttt{starter\_stack.c} to \texttt{stack.c} and fill in the eight function bodies.
  \item Sanity-check locally:
\begin{verbatim}
gcc -Wall -Wextra -o test_visible test_visible.c stack.c
./test_visible
\end{verbatim}
    All three worked examples should pass.
  \item Submit \texttt{stack.c} and \texttt{stack.h}. The autograder compiles your \texttt{stack.c} against a \emph{hidden} test suite and reports pass/fail. The hidden suite covers: the $-1$ sentinel trap; filling the stack to exactly \texttt{STACK\_MAX} and one past it; \texttt{peek} purity and \texttt{NULL}-out calls; empty/single-bracket inputs, every bracket type, $100$-deep (fits) vs.\ $101$-deep (capacity) nesting, ignored non-bracket characters, and \texttt{NULL}; single-number, negative, and multi-digit postfix; division truncation and division-by-zero; too-few-operands and leftover-operand inputs; bad tokens; empty/whitespace input; a $101$-operand overflow; and seeded random batches of bracket strings and postfix expressions checked against an oracle.
\end{enumerate}

\end{document}
